% Section 6: Limitations + Honest-Negative Results (Q2 JCIM paper v2)
% Generated from TODO-133 + TODO-138 §1.4 + TODO-141 Phase 5 + Wf-127/128/129/131/145 honest-negative ships + Sprint A/B additions
% Date: 2026-09-27
% Honest framing per party-twenty-spirit §3.5
% Honest-negative count: 8 (TODO-138 §0 baseline) + 2 (Sprint A: GroupNorm vacuous, DANN partial wiring) + 1 (Sprint B: BindingDB IC50 not vendored) + 2 (Sprint A: 7.5M deferred, theoretical tool extensions not adopted) = 13 honest-negative entries total

\section{Limitations and Honest-Negative Results}
\label{sec:limitations}

In keeping with the open-science principle that negative results should be
reported with the same prominence as positive ones, we explicitly enumerate
the honest-negative results below.  Each item cites the verbatim metric
value, the workflow that produced it, and a structural explanation of why
the metric could not be lifted.  This section is not an appendix; it is a
primary section of the paper.

The thirteen honest-negative entries decompose into four groups:
\begin{enumerate}
  \item Eight \emph{measurement} honest-negatives (§\ref{sec:limit-platinai}--§\ref{sec:limit-hilbert-original}) carried over from the TODO-138 §0 baseline.
  \item One \emph{Wf-145} honest-negative (§\ref{sec:limit-m19-rh}) for m19 LOMO Rh regression + sample-size ceiling.
  \item One \emph{baseline-alignment} honest-negative (§\ref{sec:limit-no-tier1}) for the absence of any Tier-1 metallodrug baseline.
  \item Three \emph{Sprint A} honest-negatives (§\ref{sec:limit-groupnorm-vacuous}, §\ref{sec:limit-dann-partial}, §\ref{sec:limit-train-deferred}) covering mechanism incompleteness + deferred full training.
  \item One \emph{Sprint B} honest-negative (§\ref{sec:limit-bindingdb}) for the BindingDB IC50 not vendored (4-domain 1.61\,M not 2.6\,M).
  \item One \emph{theoretical-tools} honest-negative (§\ref{sec:limit-theoretical}) for Sheaf-FRL / SGPC / TopoAD / Kan extensions not adopted (no stable reference implementation 2025--2026).
\end{enumerate}

\subsection{Architectural Signal Floor --- PlatinAI 4-Cell Combined $R^2$}
\label{sec:limit-platinai}

On the PlatinAI 4-cell synthetic pIC$_{50}$ dataset, our combined
Hilbert+ASTSR evolutionary regression achieves combined $R^2 = -5.66 \pm 1.36$
across $4$ cell lines (MCF7, A549, A2780, A2780cis; $5$ seeds $\times$
$5$ folds $\times$ $4$ cell-lines, total $20$ fits).  Three rescue
attempts with Wasserstein regularization (Wf-129), Hilbert $L_2$ +
ASTSR scale guards (Wf-127, Wf-122), and $172$-d Multi-PH + PLL feature
augmentation (Wf-128) all yielded honest-negative results ($-4.04$,
$-0.078$, $-40.24$ respectively).  Root cause: the $12$-d RDKit
descriptor signal floor --- none of our $13$ mathematical theorems
(Wasserstein gradient flow, Hilbert primal-dual, ASTSR evolutionary
search, Multi-PH Fourier, PLL signature) provides sufficient signal lift
to break the architectural ceiling.  Source:
\texttt{metrics/by\_metric/platinai\_combined\_r2.json}.

\subsection{L$_5$ Hilbert+ASTSR $R^2$ on MCF7}
\label{sec:limit-l5}

Adding $L_2$ regularization to Hilbert primal-dual and scale-clip to
ASTSR signature (Wf-127 Phase 4) prevents divergence but does not lift
PlatinAI MCF7 $R^2$ from the Wf-83 baseline of $-6.71$.  With guards,
MCF7 $R^2 = -4.04 \pm 4.31$ ($5$ seeds, $5$ folds, $10$ total fits;
Hilbert-only $-22.99 \pm 0.25$, ASTSR-only $-3.54 \pm 3.42$, combined
$\alpha=0.0$); A549 $R^2 = -1.57 \pm 1.18$.  Notably, ASTSR scale-clip
with default range $(-10, 10)$ regresses ASTSR $R^2$ on MCF7 from
$-0.014$ (no clip, Wf-86 P3 baseline) to $-3.54$ (with clip) --- the
clip forces atypical scaffold predictions to the boundary, hurting
honest fits.  Source:
\texttt{metrics/by\_metric/hilbert\_4cell\_r2.json}.

\subsection{L$_7$ Multi-PH + PLL $R^2$}
\label{sec:limit-l7}

End-to-end Multi-PH + PLL feature extraction combined with Hilbert
$L_2$ guard + ASTSR scale guard (Wf-129 Phase 3) achieves PlatinAI
4-cell combined $R^2 = -0.0780 \pm 0.0148$ ($5$ seeds $\times$ $5$
folds $\times$ $4$ cell-lines, $20$ total fits), vs the Wf-79 P2
baseline of $+0.087$ (NCI-60 real pIC$_{50}$, different cohort).  The
honest-negative result confirms that Wf-127 + Wf-122 + Wf-131 guards
prevent divergence (no $R^2 = -\infty$; all $L_2$ penalties finite
$\geq 0$; ASTSR guard never fires on widened band $[0.05, 20.0]$) but
cannot lift $R^2 \geq +0.30$ at the architectural level.  The Wf-128
$172$-d Multi-PH + PLL feature augmentation with multi-task backbone
($118{,}789$ params) is the only architectural lever that may break the
$12$-d signal floor; Wf-128 P3 29-element aggregate $R^2 = -40.24$
confirms this lever also fails.

\subsection{L$_{10}$ 3-Way $R^2$ on Real pIC$_{50}$}
\label{sec:limit-l10}

The 3-way Deflexpressor + Multi-PH + ASTSR integration on real
MetalCytoToxDB 5-metal pIC$_{50}$ (Wf-128 P3) yields honest-negative
results across $5$ precious metals (Ru, Ir, Os, Rh, Re) with $172$-d
feature stack (RDKit 12-d + Multi-PH 32-d + PLL 128-d) and multi-task
backbone ($118{,}789$ params).  On the Wf-41 P1 3-way baseline
(PlatinAI synthetic $\text{QED} \times \text{QED}$ target), $R^2 =
-0.36 \pm 0.28$; on the Wf-128 P3 29-element aggregate, $R^2 =
-40.24$ (uncomputable target signal).  Per-metal Wf-128 P2: Ru
$-0.07 \pm 0.04$ (collapse $4/5$ seeds, ASTSR $[0,1]$ clip
artefact), Ir $-201.25 \pm 80.74$, Os $-287.20 \pm 117.96$, Rh
$-183.82 \pm 97.31$, Re $-80.88 \pm 30.44$ ($3/5$ seeds collapse).
Both targets confirm the architectural signal floor --- $12$-d RDKit
descriptor cannot be lifted to $R^2 \geq 0$ by any combination of
Wasserstein + $L_2$ + scale-guard + $172$-d features + multi-task
backbone.

\subsection{m3 Vina Cisplatin --- Covalent-Binder Mismatch}
\label{sec:limit-vina-cisplatin}

Cisplatin real Vina docking sweep (Wf-121 Phase 3, exhaustiveness 8)
on $10$ pockets $\times$ $3$ seeds $\times$ $5$ poses yields mean
$-3.103 \pm 0.351$ kcal/mol (best $-3.734$ on \texttt{3ny8}, worst
$-2.548$ on \texttt{4rn0}).  Wf-141 v2 P1 exhaustiveness-$32$ re-sweep
yields $-3.1575 \pm 0.339$ kcal/mol ($\delta = -0.0545$, honest partial
lift; $0/10$ pockets reach $-5/-6/-7$ cutoff).  Wf-141 v2 P2
exhaustiveness-$64$ upper-bound (2 pockets) yields $-3.4345$, still
$3.5$ kcal/mol short of the $-7$ industry binding threshold.
Root cause: cisplatin is a covalent binder (aquation + DNA N7 guanine
crosslink); Vina's non-covalent physics model cannot represent the
covalent mechanism.  Additionally, our Vina 1.2.7 setup substitutes
\texttt{Mg} atom-type for \texttt{Pt} to bypass strict-parser rejection
(P1 SMARTS \texttt{M} rejected), introducing generic VDW
underestimation.  True cisplatin binding requires covalent docking
tools (DockCov, CovDock) and wet-lab IC$_{50}$ --- both deferred to
R23-P5 / R23-P6.  Source:
\texttt{metrics/by\_metric/vina\_kcal\_per\_mol.json}.

\subsection{Diversity Tanimoto --- Pool-Saturated Ceiling}
\label{sec:limit-diversity}

Pool-wide Morgan-ECFP4 Tanimoto distance on Lambda-MCTS generated
molecules reaches $0.137$ (EV-3, $30$ cells $\times$ $3$ seeds,
--sa-weight 0.3) vs TargetDiff $0.860$ --- a gap of $-0.723$
reflecting pool-diversity ceiling (the underlying fragment pool is
saturated).  The structurally meaningful diversity metric is
$n_{\text{distinct}}$ (count of unique generated molecules per cell),
which we lift from $1 \to 20$ via the 4-fix bundle
(Section~\ref{sec:lead-ndistinct}).  The $+0.1366$ lift over the
R12 Path-A baseline ($0.1065$) is real and reproducible
(deterministic per-cell, std $=0.0$), but further lift requires
fragment-pool expansion (TODO-25 P2.3), deferred.  Source:
\texttt{metrics/by\_metric/diversity\_tanimoto.json}.

\subsection{m15 RMSD Old Definition --- Self-Conformer Comparison}
\label{sec:limit-rmsd}

The \texttt{rigid\_fragment\_rmsd} metric ($0.428 \pm 0.547$ \AA{} across
the 10-SMILES organic pool, with per-SMILES RMSD $0.011$--$1.438$ \AA)
measures symmetry-aware MCS overlay RMSD of the maximum common
substructure after MMFF94 minimization between two
ETKDGv3+UFF-optimized conformers of the \emph{same} molecule
(TargetDiff Table 1 convention).  The metric is preserved verbatim
from TODO-22 metric \#13 (shipped 2026-09-18) and is NOT a docking
quality metric --- pocket does not enter the calculation.  To measure
true docking quality, we introduce the \texttt{docked\_pose\_rmsd}
metric (TODO-134 Phase 2, Wf-126 P1 5/5 pockets ready) --- RMSD
between Vina-docked pose and crystal ligand pose after Kabsch
alignment, which actually reflects docking accuracy.  Source:
\texttt{metrics/by\_metric/rigid\_fragment\_rmsd.json}.

\subsection{Hilbert 4-Cell $R^2$ Original Baseline}
\label{sec:limit-hilbert-original}

Original Wf-83 baseline reports Hilbert+ASTSR combined $R^2 = -6.71$
on PlatinAI 4-cell across $40$ seeds.  Wf-122 Phase 1 (Hilbert $L_2$
regularization) and Wf-127 Phase 4 ($L_2$ + ASTSR scale guard) ship as
defensive measures preventing divergence but do not lift $R^2$ from
baseline.  Per-cell-line post-fix Wf-122 P3 anchors: MCF7 $-5.946
\pm 8.193$, A549 $-4.748 \pm 7.813$, A2780, A2780cis.  Architectural
signal floor remains the binding constraint (see
Section~\ref{sec:limit-platinai}).

\subsection{Wf-145 m19 Rh Regression + Sample-Size Ceiling ($n=60$)}
\label{sec:limit-m19-rh}

Wf-131 P2 LOMO Rh fold crashed from $0.368 \pm 0.018$ (R17 3-seed
baseline) to $0.136 \pm 0.065$ (Wf-125 P2 2-seed smoke) due to a
label-bias artefact in the positive-label oversampler (the oversampled
positives were constant-inactive rather than truly active).  Wf-145 P1
fixed the cohort-generation pipeline and shipped a BALANCED Rh cohort
($500$ compounds = $250$ active + $250$ inactive, $1.0$ Catherine
6/6 PASS, scaffold\_overlap $=0$); Wf-145 P2 shipped the
\texttt{train\_metal\_lomo\_v2.py} v2 driver ($860$ LOC, $18/18$ unit
tests PASS).  However, the v2 driver \emph{full} sweep
($5$ metals $\times$ $5$ seeds $\times$ $50$ epoch = $125$ trainings)
remains DEFERRED to R20 GPU --- the current MEASURED state on disk is
\textbf{$7/125 = 5.6\%$} complete (Ru 5/5 + Ir 2/5 from Wf-131 P3
on the BUGGY cohort; Os/Rh/Re $0/5$ each).

Root cause (two compounding ceilings):
\begin{itemize}
  \item \textbf{Sample-size ceiling}: MetalCytoToxDB Rh subset has
        $n_{\text{train}}=60$ rows ($n_{\text{test}}=38$).  Even with
        perfect $5$-fold CV, the variance floor at $n=60$ saturates
        at $\sigma_{\text{AUC}} \geq 0.06$, making any lift below
        $+0.06$ undetectable.
  \item \textbf{Active-inactive separability ceiling}: per the Wf-145
        P1 BALANCED cohort diagnostic, the positive/negative label
        separability on 12-d RDKit features alone is below chance on
        the $n=60$ cohort --- no D-MPNN multi-task backbone can lift
        above $0.50$ AUC without a true wet-lab affinity signal.
\end{itemize}
Honest projection: with $50\%$ chance, Rh AUC reaches the $\geq 0.50$
floor at $5/5$ seed $\times$ $50$ epoch; with $50\%$ chance, Rh stays
below $0.50$ and we ship the honest-negative finding.  Source:
\texttt{metrics/by\_metric/metal\_lomo\_auc\_v2.json}
(\texttt{status = "DEFERRED\_PARTIAL\_RUN"}, \texttt{verdict =
"DEFERRED\_NO\_MEASURED\_VALUE"}).

\subsection{No Tier-1 Metallodrug Baseline Exists}
\label{sec:limit-no-tier1}

No Tier-1 reference implementation exists for metallodrug
structure-based drug design, de novo anticancer generation conditioned
on a specific d-block transition metal, or hybrid
metallodrug+anticancer tasks.  We surveyed \texttt{DrugOOD}
(ecoli/scaffold/assay OOD splits), \texttt{NCI-60} GI50 panel,
\texttt{TargetDiff} (ICLR 2023, CrossDocked), \texttt{DiffDock}
(NeurIPS 2024), \texttt{Pocket2Mol} (ICML 2022), \texttt{DiffSBDD}
(ICML 2023), \texttt{Krasnov}~2026 (chemprop v2 NCI-60 9-cell
baseline $0.81$ AUC), and \texttt{iMoLD}~(NeurIPS 2022, DrugOOD
assay $0.7211$ / scaffold $0.6884$ / size $0.6792$ OOD AUC).  None of
these baselines is metal-conditioned; they are cite-only context
numbers (paper Table~1 footnote, not head-to-head).  iMoLD
GitHub repo (\texttt{tencent-ailab/DrugOOD}) was searched for a
graph-NN multi-task drop-in; no public re-implementation exists on
this host, and a clone was attempted without success (GitHub search
exhausted; fallback to paper reference numbers + cite-only).
\textbf{Honest implication}: paper §4.1 cannot claim
``beats iMoLD'' as a head-to-head number; we cite the published
$0.7211$ / $0.6884$ / $0.6792$ numbers as Tier-1 context, not as a
direct comparator.  Same-protocol head-to-head requires a graph-NN
backbone + full training, deferred to R20 GPU.  Source:
\texttt{metrics/by\_metric/multi\_task\_4domain\_v2.json}
\texttt{valid\_comparisons\_cite\_only}.

\subsection{Sprint A: GroupNorm Replacement Vacuous}
\label{sec:limit-groupnorm-vacuous}

The Sprint A GroupNorm drop-in
(\texttt{molmetal/molmetal\_lam/sbdd\_env/groupnorm\_replacement.py},
$140$ LOC, $12/12$ tests PASS in
\texttt{molmetal/tests/test\_groupnorm\_replacement.py}) ships a
recursive replacement helper that swaps \texttt{nn.BatchNorm1d} /
\texttt{nn.BatchNorm2d} for \texttt{nn.GroupNorm(num\_groups, num\_channels)}.
However, \texttt{grep -rn "BatchNorm" molmetal/models/ molmetal/baselines/}
returns \textbf{zero hits} --- our in-tree D-MPNN backbone is
\emph{no-norm by design} (it ships without any batch-normalisation
layers; chemprop v2 also confirms no GroupNorm/LayerNorm out-of-box,
BatchNorm1d gated by \texttt{batch\_norm=False} default).  The GroupNorm
helper is correct and ready to consume any imported chemprop-style
MPNN backbone, but applying it to the current Mol-Metal D-MPNN is
\textbf{a no-op}.  This is the documented ``vacuous'' case in
\texttt{party-twenty-spirit} §3.5: the mechanism is correct, the hook
is wired in \texttt{train\_dmpnn\_multitask.py}, but the deliverable
is currently a no-op.  Source:
\texttt{TODO/completed/q2\_sprint\_self\_check.md} §1.

\subsection{Sprint A: DANN Per-Step Wiring Deferred}
\label{sec:limit-dann-partial}

The Sprint A DANN (Domain-Adversarial Neural Network) drop-in ships
the gradient-reversal layer
(\texttt{molmetal/molmetal\_lam/sbdd\_env/gradient\_reversal.py},
\texttt{ReverseLayerF}, $83$ LOC, $20/20$ tests PASS in
\texttt{molmetal/tests/test\_gradient\_reversal.py}), the
domain-classifier head
(\texttt{molmetal/models/dmpnn.py::DomainAdversarialHead}, lines
370--420, instantiable with \texttt{in\_dim, num\_domains=4}), and the
CLI flags \texttt{--use-dann}, \texttt{--lambda-grl} (default
$\lambda_{\text{grl}}=0.1$) in
\texttt{molmetal/scripts/train\_dmpnn\_multitask.py}.  However, the
\emph{per-step} training-loop wiring that computes
$\mathcal{L}_{\text{domain}} = \text{CE}(\text{domain\_logits},
\text{domain\_labels})$ and adds $\lambda \cdot \mathcal{L}_{\text{domain}}$
to $\mathcal{L}_{\text{task}}$ is \textbf{NOT yet implemented}.  The CLI
logs ``DANN enabled: lambda\_grl=0.1'' but does not invoke
\texttt{DomainAdversarialHead.forward()} inside the inner training
loop.  Per \texttt{party-twenty-spirit} §3.5, this is the ``partial
wiring'' case: components ship + tests pass, but the head is not
invoked per step.  The follow-up is $\sim 30$ LOC deferred to R20 GPU
(the same blocker that gates the LOMO $5 \times 5 \times 50$ sweep).
Source: \texttt{TODO/completed/q2\_sprint\_self\_check.md} §4.

\subsection{Sprint A: 7.5\,M Training Deferred to R20 (current 13\,K = 0.17\,\%)}
\label{sec:limit-train-deferred}

The cross-domain training recommendation (TODO-142 §0) identifies that
the in-tree D-MPNN has been trained on $13{,}035$ unique molecules
across $7$ data sources, vs an available $7{,}500{,}000$ molecules
($0.17\%$ coverage).  Per-domain current usage vs available:

\begin{itemize}
  \item MetalCytoToxDB ($5$ metal): \textbf{5{,}207} trained / $5{,}207$ available = $100\%$.
  \item NCI-60 $9$-cell subset: \textbf{7{,}824} trained / $7{,}824$ available = $100\%$.
  \item PlatinAI $4$-cell: $\sim$\textbf{5{,}000} trained / $5{,}000$ available = $100\%$.
  \item DrugOOD native: \textbf{0} trained / $1{,}580{,}000$ available = $0\%$.
  \item NCI-60 GI50 raw: \textbf{0} trained / $4{,}850{,}000$ available = $0\%$.
  \item BindingDB IC50: \textbf{0} trained / $1{,}072{,}000$ available = $0\%$ (see §\ref{sec:limit-bindingdb}).
  \item CrossDocked (SBDD targets): \textbf{0} trained / $183{,}000$ PDB = $0\%$.
\end{itemize}

This is the root cause of every paper §4 lead-claim incompleteness ---
not a model defect but a data-coverage gap.  The full
$5$-metal $\times$ $5$-source $\times$ $7.5$\,M training requires
weeks of GPU-0 Blackwell time and is DEFERRED to R20 GPU (2027).
Source: \texttt{TODO/pending/142\_cross\_domain\_4\_5m\_training\_recommendation.md} §0.

\subsection{Sprint B: BindingDB IC50 Not Vendored --- 4-Domain 1.61\,M (not 2.6\,M)}
\label{sec:limit-bindingdb}

The Sprint A $4$-domain \texttt{WeightedRandomSampler} drop-in
(\texttt{molmetal/data/four\_domain\_combined.py} +
\texttt{molmetal/data/four\_domain\_weights.py}, $26/26$ tests PASS in
\texttt{molmetal/tests/test\_four\_domain\_sampler.py}) originally
targeted a $2.6$\,M-row $4$-domain combine ($1.58$\,M DrugOOD + $1.07$\,M
BindingDB + $7.8$\,K NCI-60 + $5$\,K MetalCytotox).  In practice,
\textbf{true BindingDB IC50 ($1.07$\,M rows) is not vendored locally};
the closest analog is the protein-metal-binding CSV ($5{,}258$ rows)
shipped by TODO-110.  The actual locally-vendored $4$-domain total
is $1{,}589{,}302 + 5{,}258 + 7{,}824 + 5{,}207 \approx 1.61$\,M, not $2.6$\,M.
This is pinned by \texttt{test\_combined\_total\_is\_2\_6m} which
intentionally checks the $1.59$\,M--$1.62$\,M band and documents the
BindingDB fallback.  Source:
\texttt{TODO/completed/q2\_sprint\_self\_check.md} §3.

Additional caveat on the $\alpha = 0.5$ soft-balancing exponent:
with DrugOOD at $1.58$\,M and MetalCytoToxDB at $5$\,K ($316 \times$
ratio), the per-row sampling weight for MetalCytoToxDB rows is
$\sqrt{316} \approx 17.8 \times$ higher than DrugOOD rows --- the
right direction (rare-domain upweight) but not strict per-batch
balance.  Strict per-batch balance would require $\alpha \to 0$
(inverse-frequency) which would massively over-sample rare domains.
The $\alpha = 0.5$ choice is a practical compromise that has not been
grid-searched; deferred to R20.

\subsection{Sprint B: Sheaf-FRL / SGPC / TopoAD / Kan Extensions Not Adopted}
\label{sec:limit-theoretical}

The TODO-143 §1.2 theoretical-tool survey identified four 2025--2026
methods for topological + categorical structure-aware multi-task
learning:

\begin{itemize}
  \item \textbf{Sheaf-FRL} (SPAICOM 2026, federated sheaf Laplacian
        + Procrustes closed-form): cloned at
        \texttt{.cloned/Sheaf-FRL/} but only the FRL sheaf-Laplacian
        primitives are reusable; the federated-learning scaffolding is
        out-of-scope for our single-host multi-task D-MPNN.  Full lift
        estimated at $5$--$7$ days with high architectural-mismatch
        risk; deferred to R21 theoretical upgrade.
  \item \textbf{SGPC} (AAAI 2026, structured Gaussian process
        classifier): no stable reference implementation found on GitHub
        or PyPI; cannot clone.
  \item \textbf{TopoAD / pTU} (2025, topology-aware anomaly detection
        on persistent-homology features): no stable reference
        implementation found; cannot clone.
  \item \textbf{Kan extensions} (2026 categorical-ML line): no
        implementation found beyond theory papers; cannot clone.
\end{itemize}

The honest finding is that \emph{the theory is published but no stable
code exists that we can drop in}.  Multiflow's
\texttt{contact\_sheaf.py} lifted only the \emph{shape-port} of the
ContactSection dataclass (\texttt{molmetal/molmetal\_lam/sbdd\_env/sheaf\_section.py},
$215$ LOC, $16/16$ tests PASS in
\texttt{molmetal/tests/test\_sheaf\_section.py}) --- the full
sheaf-Laplacian / Procrustes machinery is NOT lifted.  Source:
\texttt{TODO/completed/q2\_sprint\_self\_check.md} §9.3, §6 +
\texttt{TODO/pending/143\_clone\_adapt\_cross\_domain\_training.md} §1.2.

\subsection{Wet-Lab Validation Gap}
\label{sec:limit-wetlab}

No wet-lab IC$_{50}$, protein-metal binding affinity, or
cellular-uptake measurement is reported in this paper.  All
in-silico metrics are computational proxies.  The
ORNL/PKU/Xiehe-NCI60 collaboration pipeline (Wf-74 v2 + Wf-116) is
in flight and will provide wet-lab validation in R23-P6 (deferred).
Without wet-lab validation, the clinical relevance of the in-silico
lead claims is unverified.

\subsection{Future Architectural Directions}
\label{sec:limit-future}

To break the architectural signal floor and the covalent-binder Vina
limit, future work will integrate:
\begin{itemize}
  \item AlphaFold-Multimer protein-aware descriptors (R22-P1)
  \item Covalent docking tools DockCov / CovDock (R23-P5)
  \item Wet-lab IC$_{50}$ + protein-metal binding assays (R23-P6)
  \item Foundation-model baseline (MolGen / DrugGPT) for SOTA comparator
  \item Larger tmQM pretraining ($100$-epoch CFM with
        Wasserstein+ASTSR guard)
  \item D-MPNN per-step DANN loss wiring ($\sim 30$ LOC, R20 GPU)
  \item Grid-searched $\alpha \in \{0.0, 0.25, 0.5, 0.75, 1.0\}$ for
        $4$-domain WeightedRandomSampler balance
  \item Wf-145 P2 v2 LOMO driver $5/5$ metal $\times$ $5$ seed
        $\times$ $50$ epoch full sweep (R20 GPU, $\sim 4$\,h wall)
  \item $5$-metal $\times$ $5$-source $\times$ $7.5$\,M full training
        (R20 GPU weeks, post Phase 1 mechanism verification)
\end{itemize}